%! Unit 2, Lessons 2.1–2.5 — Unit Test Study Guide (reference sheet; no answer key by design)
% This is a REFERENCE sheet, not a worksheet: every definition, every worked move, and every
% trap is already filled in, so there is nothing to answer and therefore no _key sibling.
% The practice a student does before the test is ../practice_test/, which HAS a key.
% NAMESTRIP: no \namedateperiod — the name row lives on a packet cover and the unit tests only.
% Vocabulary below is quoted from the five lessons' own vocabulary boxes (notes_key \termans
% entries) so a student reading this reads the same words twice.
% Worked contexts (museum audio tour; hiking trails) are deliberately NOT used in any lesson,
% AP practice, or homework — and the practice and actual tests must not use them either.
\documentclass[12pt]{article}
\usepackage{apstats-article}
\usepackage{apstats-boxes}
\usepackage{ltablex}
\keepXColumns
% One tabularx per group of lessons rather than one long table: a tabularx with X
% columns is an unbreakable box, so a single long table floats whole to a fresh page
% and strands most of page 1. Short tables flow, and each carries its lesson tag.

% Optional arg: the \boxguard size (lines kept with the strip). The last section is three
% lines long, so it guards only what it needs.
\newcommand{\sghead}[2][10]{%
  \vspace{0.13in}
  \boxguard[#1]
  \begin{headlinebox}{navy}\color{white}\bfseries\small #2\end{headlinebox}
  \vspace{0.09in}}

\begin{document}

\pageheader{Unit 2 \quad Lessons 2.1--2.5}{Unit Test Study Guide}

\vspace{0.12in}

\boxguard[8]
\begin{remindbox}
\small
The test covers \textbf{Lesson 2.1} through \textbf{Lesson 2.5} --- two categorical variables,
then scatterplots, correlation, regression lines, residuals, and when a line is the wrong
model. Multiple choice, then two free-response questions: one on a two-way table, one on a
scatterplot with computer output and a residual plot. Every graph is already drawn: \textbf{read}
it, and answer \textbf{in context} --- variables, units, and \emph{predicted} for the line.
\end{remindbox}

\sghead{1. Vocabulary you are expected to use, not just recognize}

\small
\renewcommand{\arraystretch}{1.15}

\textbf{From Lesson 2.1 --- two categorical variables}
\par\vspace{3pt}
\noindent\begin{tabularx}{\linewidth}{@{} l X @{}}
\rowcolor{sky}
\textbf{Term} & \textbf{What it means} \\
\midrule
Two-way table & a table of counts for two categorical variables on the same individuals. \\
Joint rel.\ freq. & one cell $\div$ the \textbf{grand} total. \\
Marginal rel.\ freq. & a row or column total $\div$ the \textbf{grand} total. \\
Conditional rel.\ freq. & a cell count $\div$ the total of its \textbf{own} row or column.
                   In ``of those who\ldots'', the words after \emph{of} name the denominator. \\
Segmented bar graph & one bar per group, totalling 100\%, split into its conditional distribution. \\
Association & the distribution of one variable changes across the categories of the other. \\
\end{tabularx}

\vspace{7pt}
\boxguard[9]
\textbf{From Lessons 2.2 and 2.3 --- scatterplots, $r$, and the line}
\par\vspace{3pt}
\noindent\begin{tabularx}{\linewidth}{@{} l X @{}}
\rowcolor{sky}
\textbf{Term} & \textbf{What it means} \\
\midrule
Scatterplot & two quantitative variables on the same individuals; one dot per individual. \\
Explanatory \& response & the explanatory variable ($x$-axis) helps explain or predict the
                   response ($y$-axis), the outcome. \\
DOFS & describe a scatterplot by \textbf{D}irection, \textbf{O}utliers \& unusual features,
                   \textbf{F}orm, \textbf{S}trength --- always in context. \\
Correlation $r$ & a number from $-1$ to $1$ giving the direction and strength of a
                   \textbf{linear} association. \\
Regression line & a line $\hat{y}=a+bx$ that uses the explanatory variable $x$ to predict the
                   response variable $y$. \\
Predicted value ($\hat{y}$) & the $y$-value the line gives for a particular $x$; the actual
                   $y$ usually differs from it. \\
Extrapolation & predicting for an $x$ outside the $x$-values used to build the line --- unreliable. \\
Residual & actual minus predicted, $y-\hat{y}$; positive means the point lies above the line. \\
\end{tabularx}

\vspace{7pt}
\boxguard[9]
\textbf{From Lesson 2.4 --- the least-squares line}
\par\vspace{3pt}
\noindent\begin{tabularx}{\linewidth}{@{} l X @{}}
\rowcolor{sky}
\textbf{Term} & \textbf{What it means} \\
\midrule
Least-squares line & the line $\hat y = a + bx$ that makes the sum of the squared residuals as
                   small as possible. Passes through $(\bar x, \bar y)$. \\
Slope $b$ & $b = r\,s_y/s_x$: change in the \emph{predicted} $y$ per 1-unit increase
                   in $x$; same sign as $r$. \\
Intercept $a$ & $a = \bar y - b\,\bar x$: predicted $y$ at $x = 0$ --- meaningful only if
                   $x = 0$ makes sense \emph{and} is near the data. \\
$r^2$ (R-Sq) & the proportion of the variation in $y$ explained by the linear model with $x$. \\
\end{tabularx}

\vspace{7pt}
\boxguard[9]
\textbf{From Lesson 2.5 --- when a line is the wrong model}
\par\vspace{3pt}
\noindent\begin{tabularx}{\linewidth}{@{} l X @{}}
\rowcolor{sky}
\textbf{Term} & \textbf{What it means} \\
\midrule
Residual plot & residuals graphed against $x$. Random scatter around 0: a line fits.
                   A curve: wrong model. \\
Outlier & a point that does not follow the trend of the rest; it has a large residual. \\
High-leverage point & a point whose $x$-value is much larger or smaller than the other
                   $x$-values. \\
Influential point & a point that, if removed, changes the slope, $y$-intercept, or $r$ a lot. \\
Transformation & re-expressing a variable (for example, $\log y$) to straighten a curve; undo it
                   at the end ($10^{\,\widehat{\log y}}$). \\
\end{tabularx}

\normalsize

\sghead{2. The moves the test will ask you to make}

\small
\textbf{Move 1 --- Use the right denominator (2.1).} 250 museum visitors: 80 took the audio
tour and 52 of them rated the visit \emph{excellent}; 170 skipped it and 68 of them did.
\emph{Worked:} \textbf{joint} --- toured and excellent: $52/250 = 0.208$. \textbf{Marginal}
--- excellent: $120/250 = 0.48$. \textbf{Conditional} --- of the tour-takers, excellent:
$52/80 = 0.65$; of the others, $68/170 = 0.40$.

\vspace{3pt}
\textbf{Move 2 --- Decide association (2.1).} Compare the \emph{conditional} percents, never
the counts. \emph{Worked:} 68 excellent ratings is the bigger \emph{count}, but 65\% vs.\
40\% is the fair comparison, so tour and rating \textbf{are associated} in this sample --- and
because visitors chose whether to take the tour, it does not show the tour \emph{causes}
better ratings.

\vspace{3pt}
\textbf{Move 3 --- Build the line from summary statistics (2.4).} 15 hiking trails: elevation
gain $x$ (hundreds of feet, 2 to 20) and hiking time $y$ (minutes).
$\bar x = 10$, $s_x = 5$, $\bar y = 150$, $s_y = 40$, $r = 0.90$.
\emph{Worked:} $b = 0.90 \cdot 40/5 = 7.2$; \ $a = 150 - 7.2(10) = 78$; \ so
$\widehat{\text{time}} = 78 + 7.2(\text{gain})$.

\vspace{3pt}
\textbf{Move 4 --- Read computer output (2.4).} The same line, as software prints it:
\par\vspace{3pt}
{\centering\footnotesize\ttfamily
\begin{tabular}{@{} l r r r r @{}}
Predictor & Coef & SE Coef & T & P \\ \hline
Constant & 78.00 & 10.74 & 7.26 & 0.000 \\
Gain & 7.200 & 0.967 & 7.45 & 0.000 \\
\end{tabular}\par
S = 18.09 \quad R-Sq = 81.0\% \quad R-Sq(adj) = 79.5\%\par}
\vspace{3pt}
$a$ is the Coef in the \textbf{Constant} row; $b$ is the Coef in the row named for $x$. R-Sq
is $r^2$. To get $r$, take $\sqrt{0.81} = 0.90$ and give it the \textbf{sign of the slope}.

\vspace{3pt}
\textbf{Move 5 --- Predict, and check the range (2.3).} \emph{Worked:} a 1{,}200-ft trail
($x = 12$): $\hat y = 78 + 7.2(12) = 164.4$ minutes. A 4{,}000-ft trail ($x = 40$) gives
366 minutes --- but 40 is far outside 2 to 20, so that is \textbf{extrapolation}: don't trust it.

\vspace{3pt}
\textbf{Move 6 --- Residual, with the sign read correctly (2.3).} \emph{Worked:} that
1{,}200-ft trail actually took 150 minutes. Residual $= y - \hat y = 150 - 164.4 = -14.4$
minutes: 14.4 minutes \emph{less} than predicted --- the point is below the line, so the line
\textbf{over}predicted.

\vspace{3pt}
\textbf{Move 7 --- Judge the model with the residual plot (2.5).} Look for a \emph{pattern}:
\par\vspace{4pt}
{\centering
\begin{tikzpicture}[x=0.26cm, y=0.2cm]
  \foreach \dx/\lab in {0/{Random scatter: a line fits}, 17/{A curve: wrong model},
                        34/{One big residual: an outlier}} {
    \begin{scope}[xshift=\dx*0.26cm]
      \draw[thick] (0,-3) -- (0,3); \draw[thick] (0,-3) -- (13,-3);
      \draw[dashed, slate] (0,0) -- (13,0);
      \node[below, font=\scriptsize, align=center] at (6.5,-3.1) {\lab};
    \end{scope}}
  \foreach \x/\e in {1/1.2, 2/-1.4, 3/0.5, 4/-0.6, 5/1.6, 6/-1.1, 7/0.3, 8/-1.7, 9/1.0,
                     10/-0.4, 11/1.3, 12/-0.9}
    \fill[navy] (\x,\e) circle (4pt);
  \foreach \x/\e in {1/2.3, 2/1.1, 3/0.1, 4/-0.8, 5/-1.4, 6/-1.7, 7/-1.6, 8/-1.2, 9/-0.5,
                     10/0.4, 11/1.4, 12/2.5}
    \fill[navy] (17+\x,\e) circle (4pt);
  \foreach \x/\e in {1/0.6, 2/-0.7, 3/0.4, 4/-0.3, 5/-0.8, 7/0.5, 8/-0.5, 9/0.7,
                     10/-0.6, 11/0.2, 12/-0.4}
    \fill[navy] (34+\x,\e) circle (4pt);
  \fill[redacc] (34+6,2.6) circle (4pt);
\end{tikzpicture}\par}
\vspace{3pt}
To call a point \textbf{influential}, quote the numbers: the slope, intercept, or $r$ with
and without it.
\normalsize

\vspace{0.10in}
\boxguard[20]
\begin{scenariobox}[The sentences that earn full credit]{navy}
\small
\textbf{Describe a scatterplot:} ``There is a \textbf{strong, positive, linear} association
between elevation gain and hiking time: trails with more climbing tend to take longer. No
point stands far from the pattern.''\par\vspace{2pt}
\textbf{Slope:} ``For each additional 100 feet of elevation gain, the \emph{predicted} hiking
time rises 7.2 minutes.''\par\vspace{3pt}
\textbf{Intercept:} ``A trail with no elevation gain has a predicted time of 78 minutes'' ---
then say whether $x = 0$ is sensible and near the data (here it is below 2 to 20: caution).\par\vspace{3pt}
\textbf{$r^2$:} ``About \textbf{81\%} of the variation in hiking time is explained by the
linear model with elevation gain.''\par\vspace{3pt}
\textbf{Residual:} ``This trail took 14.4 minutes \emph{less} than the line predicted.''
\par\vspace{3pt}
\textbf{Association (two-way table):} ``In this sample, 65\% of audio-tour visitors rated the
visit excellent, compared with 40\% of the others, so the tour and the rating are associated.
The visitors chose, so this does not show the tour \emph{causes} higher ratings.''
\par\vspace{3pt}
Each names both variables and the units --- delete the context and it is no longer an answer.
\end{scenariobox}

\sghead{3. Seven traps that cost points every year}

\small
\begin{enumerate}[leftmargin=1.6em, itemsep=2pt]
  \item \textbf{Counts are not rates.} Compare percents, each out of its \emph{own} group's
        total --- not the grand total. The bigger count can be the smaller percent.
  \item \textbf{$r \approx 0$ is not ``no relationship.''} It means no \emph{linear} one; a
        strong curve can have $r \approx 0$.
  \item \textbf{$r$ has no units}: new units or swapped axes leave it unchanged. And $-0.9$ is
        \emph{stronger} than $0.6$.
  \item \textbf{The residual is $y - \hat y$}, never $\hat y - y$. Positive: the point is
        \emph{above} the line, so the line predicted too \emph{low}.
  \item \textbf{Say \emph{predicted}.} The slope describes the line, not each individual; and
        observational data show no cause.
  \item \textbf{$r = \pm\sqrt{r^2}$, with the slope's sign.} $r^2$ is variation
        \emph{explained}, not a share of the points.
  \item \textbf{A high $r$ does not make a line appropriate} --- only the residual plot checks
        form. And a point off the trend is not automatically influential; leverage comes from
        being far out in $x$.
\end{enumerate}
\normalsize

\sghead[6]{4. Before you walk in}

\small
Work the \textbf{practice test} without notes first, then check it against its key. A
free-response answer shorter than the key's is almost always missing the context or the word
\emph{predicted}. Re-read each lesson cover's \textbf{Keep in Mind} box.
\normalsize

\end{document}
